---
format: yohaku/0.2
target: text
macro: '[[macros/raw]]'
---

## 正文
````text
"""Adaption Responses transport with local Yohaku function execution.

This development module is installed only into a separately named candidate.
Network POSTs are never retried. Journals reserve identities before network/tool
execution, so a crash is an unknown outcome rather than permission to replay.
"""
import argparse
import contextlib
import hashlib
import http.client
import json
import os
from pathlib import Path
import re
import signal
import stat
import sys
import threading
import time
import unicodedata
import urllib.parse
import uuid

VERSION = 'yohaku-adaption-responses/2'
API_BASE = 'https://api.prod.adaptionlabs.ai/api/v1'
MAX_RESPONSE = 32 * 1024 * 1024
SENSITIVE = {'authorization', 'api_key', 'adaption_api_key', 'access_token', 'refresh_token', 'password'}


def sha(value):
    raw = value if isinstance(value, bytes) else json.dumps(value, sort_keys=True, ensure_ascii=False).encode()
    return hashlib.sha256(raw).hexdigest()


def redact(value):
    if isinstance(value, dict):
        return {k: '[REDACTED]' if k.lower() in SENSITIVE else redact(v) for k, v in value.items()}
    if isinstance(value, list):
        return [redact(v) for v in value]
    return value


def credentials():
    key = os.environ.get('ADAPTION_API_KEY')
    org = os.environ.get('ADAPTION_ORGANIZATION_ID')
    source = 'environment' if key else None
    file = os.environ.get('YOHAKU_ADAPTION_CREDENTIAL_FILE')
    if not key and file:
        path = Path(file)
        info = path.lstat()
        if not stat.S_ISREG(info.st_mode) or path.is_symlink() or info.st_mode & 0o077:
            raise ValueError('Adaption credential file must be a private regular file (0600)')
        value = json.loads(path.read_bytes())
        key = value.get('api_key') or value.get('ADAPTION_API_KEY')
        org = value.get('organization_id') or value.get('ADAPTION_ORGANIZATION_ID')
        source = 'private_file'
    if not isinstance(key, str) or not key or any(c.isspace() for c in key):
        raise ValueError('Missing ADAPTION_API_KEY or YOHAKU_ADAPTION_CREDENTIAL_FILE api_key')
    if org is not None and (not isinstance(org, str) or '\r' in org or '\n' in org):
        raise ValueError('Invalid Adaption organization id')
    return key, org, source


class OutcomeUnknown(RuntimeError):
    pass


class LocalCancelled(RuntimeError):
    pass


class BudgetExceeded(RuntimeError):
    pass


class Journal:
    def __init__(self, directory):
        self.root = Path(directory)
        self.root.mkdir(parents=True, exist_ok=True, mode=0o700)
        if self.root.is_symlink():
            raise ValueError('Journal directory cannot be a symlink')

    def path(self, identity):
        if not re.fullmatch(r'[A-Za-z0-9_.:-]{1,180}', identity):
            raise ValueError('Use an immutable ASCII request identity')
        return self.root / (identity + '.json')

    def load(self, identity):
        p = self.path(identity)
        return json.loads(p.read_bytes()) if p.exists() else None

    def save(self, identity, value):
        p = self.path(identity)
        tmp = p.with_suffix('.pending-' + uuid.uuid4().hex)
        with tmp.open('x', encoding='utf-8') as stream:
            os.chmod(tmp, 0o600)
            json.dump(redact(value), stream, ensure_ascii=False, indent=2)
            stream.write('\n'); stream.flush(); os.fsync(stream.fileno())
        os.replace(tmp, p)

    @contextlib.contextmanager
    def lock(self, identity):
        import fcntl
        with self.path(identity).with_suffix('.lock').open('a') as stream:
            fcntl.flock(stream.fileno(), fcntl.LOCK_EX)
            yield


HTTP_DIAGNOSTICS_VERSION = 'http-response-observability/1'
HTTP_SAFE_BODY_CHAR_LIMIT = 65536


def _transport_safe_text(value, key, organization_id, limit=None):
    text = str(value)
    # Replace complete secrets before truncation, so the cutoff cannot expose
    # a prefix of the current credential or organization identifier.
    for secret in (key, organization_id):
        if not isinstance(secret, str) or not secret:
            continue
        variants = {secret, json.dumps(secret, ensure_ascii=True)[1:-1],
                    json.dumps(secret, ensure_ascii=False)[1:-1],
                    urllib.parse.quote(secret, safe=''), urllib.parse.quote_plus(secret)}
        for variant in sorted(variants, key=len, reverse=True):
            if variant:
                text = text.replace(variant, '[REDACTED]')
    text = re.sub(r'(?i)(bearer\s+)[A-Za-z0-9._~+/=-]+', r'\1[REDACTED]', text)
    # HTTP errors occasionally echo request credentials under named fields.
    text = re.sub(r"""(?ix)(["']?(?:authorization|api[_-]?key|access[_-]?token|refresh[_-]?token|password|organization[_-]?id)["']?\s*[:=]\s*)["']?[^\s,"'<>}]+""",
                  r'\1[REDACTED]', text)
    return text if limit is None else text[:limit]


def _transport_body_diagnostic(raw, key, organization_id, complete):
    raw_hash = hashlib.sha256(raw).hexdigest()  # Before any redaction.
    decoded = raw.decode('utf-8', errors='replace')
    stripped = decoded.lstrip()
    if not raw:
        classification = 'empty'
    elif stripped.lower().startswith(('<!doctype html', '<html', '<head', '<body')):
        classification = 'html'
    elif stripped.startswith(('{', '[')):
        classification = 'json_like'
    elif stripped.startswith('data:'):
        classification = 'event_stream_like'
    else:
        classification = 'text_or_binary'
    safe = _transport_safe_text(decoded, key, organization_id)
    return {'bytes_observed': len(raw), 'raw_sha256': raw_hash,
            'raw_sha256_covers_complete_body': bool(complete),
            'body_read_complete': bool(complete), 'classification': classification,
            'safe_body_copy': safe[:HTTP_SAFE_BODY_CHAR_LIMIT],
            'safe_body_copy_truncated': len(safe) > HTTP_SAFE_BODY_CHAR_LIMIT,
            'safe_body_char_limit': HTTP_SAFE_BODY_CHAR_LIMIT,
            'redaction': 'Current key/org literals, JSON/URL forms and named credential patterns; raw bytes never persisted.'}



SSE_TRANSPORT_VERSION = 'responses-sse-terminal-transport/3'
MODEL_NAMESPACE_REGISTRY_SHA256 = 'f6b95428099977bb34221735b6af19f0770bf2244908d213aa70994a4014ff9c'
MODEL_NAMESPACE_REGISTRY_FILE = Path(__file__).resolve().parent / 'MODEL-NAMESPACE-MAPPING.json'


class ModelNamespaceIdentityError(RuntimeError):
    pass


def _transport_namespace_registration(requested_public_model):
    path = MODEL_NAMESPACE_REGISTRY_FILE
    if path.is_symlink() or not path.is_file():
        raise ModelNamespaceIdentityError('Pinned model namespace registry missing')
    raw = path.read_bytes()
    if hashlib.sha256(raw).hexdigest() != MODEL_NAMESPACE_REGISTRY_SHA256:
        raise ModelNamespaceIdentityError('Pinned model namespace registry changed')
    registry = json.loads(raw)
    bindings = registry.get('bindings')
    if not isinstance(bindings, list) or len(bindings) != 1:
        raise ModelNamespaceIdentityError('Exactly one pre-registered deployment required')
    binding = bindings[0]
    if (binding.get('provider'), binding.get('canonical_public_model'), binding.get('supplier_model_verbatim')) != ('adaption', 'glm-53-flash', 'glm-5.3-flash'):
        raise ModelNamespaceIdentityError('Unsupported namespace registration')
    if requested_public_model != binding['canonical_public_model']:
        raise ModelNamespaceIdentityError('Requested public deployment has no registered mapping')
    sources = binding.get('sources', [])
    if len(sources) != 3 or {x.get('kind') for x in sources} != {'catalog', 'official_model_page', 'observed_terminal'}:
        raise ModelNamespaceIdentityError('All three namespace evidence sources required')
    evidence = {}
    for source in sources:
        p = Path(source['path'])
        if not p.is_absolute():
            p = path.parent / p
        if p.is_symlink() or not p.is_file() or hashlib.sha256(p.read_bytes()).hexdigest() != source['sha256']:
            raise ModelNamespaceIdentityError('Namespace mapping source missing or changed')
        evidence[source['kind']] = p
    catalog = json.loads(evidence['catalog'].read_bytes())
    entries = (catalog.get('response') or {}).get('data', [])
    if catalog.get('method') != 'GET' or catalog.get('http_status') != 200 or not any(x.get('id') == 'glm-53-flash' and x.get('display_name') == 'GLM 5.3 Flash' for x in entries):
        raise ModelNamespaceIdentityError('Catalog does not establish the registered public deployment')
    document = evidence['official_model_page'].read_text()
    if 'glm-53-flash' not in document or 'https://huggingface.co/zai-org/GLM-5.3-Flash' not in document:
        raise ModelNamespaceIdentityError('Official deployment/upstream card evidence missing')
    prior = json.loads(evidence['observed_terminal'].read_bytes())
    observed = (prior.get('sse_diagnostic') or {}).get('terminal_response_observed') or {}
    fields = observed.get('public_model_fields', [])
    if (prior.get('request') or {}).get('model') != 'glm-53-flash' or observed.get('status') != 'completed' or observed.get('model') != 'glm-5.3-flash':
        raise ModelNamespaceIdentityError('Registered terminal observation does not match both namespaces')
    if {'path': '$.model', 'value': 'glm-53-flash'} not in fields or {'path': '$.response.model', 'value': 'glm-5.3-flash'} not in fields:
        raise ModelNamespaceIdentityError('Both exact public/upstream terminal fields required')
    return registry, binding


def _transport_verified_exact_identity(payload, candidate, origin, event_public_model=None):
    requested = payload.get('model')
    if not isinstance(candidate, dict) or candidate.get('object') != 'response' or candidate.get('model') != requested:
        raise ModelNamespaceIdentityError('Raw complete response model is not the exact requested public ID')
    if candidate.get('status') not in ('completed', 'incomplete', 'failed') or not isinstance(candidate.get('output'), list) or not isinstance(candidate.get('id'), str) or not candidate['id']:
        raise ModelNamespaceIdentityError('Complete exact-model response shape required')
    if event_public_model is not None and event_public_model != requested:
        raise ModelNamespaceIdentityError('An outer model field contradicts the exact raw response model')
    return {'format': 'verified-model-namespace-identity/1', 'provider': 'adaption',
        'verification_source': 'host_verified_complete_exact_response',
        'verification_method': 'exact_raw_response_model', 'proof_origin': origin,
        'canonical_public_model': requested, 'event_public_model': event_public_model,
        'supplier_model_verbatim': candidate['model'], 'upstream_model': None,
        'response_id': candidate['id'], 'raw_response_object_sha256': sha(candidate),
        'mapping_registry_id': None, 'mapping_registry_sha256': None, 'mapping_evidence': [],
        'no_model_field_rewrite': True}


def _transport_verified_namespace_identity(payload, event, candidate, observed_response_id):
    if not isinstance(candidate, dict) or candidate.get('id') != observed_response_id:
        raise ModelNamespaceIdentityError('Final response identity differs across the SSE stream')
    if candidate.get('model') == payload.get('model'):
        return _transport_verified_exact_identity(payload, candidate, 'client_complete_sse_terminal', event.get('model'))
    registry, binding = _transport_namespace_registration(payload.get('model'))
    if event.get('model') != payload.get('model'):
        raise ModelNamespaceIdentityError('Final event public model differs from requested public deployment')
    if not isinstance(candidate, dict) or candidate.get('model') != binding['supplier_model_verbatim']:
        raise ModelNamespaceIdentityError('Final nested supplier model differs from pre-registered upstream label')
    if not isinstance(candidate.get('id'), str) or not candidate['id'] or candidate['id'] != observed_response_id:
        raise ModelNamespaceIdentityError('Final response identity differs across the SSE stream')
    return {'format': 'verified-model-namespace-identity/1', 'provider': 'adaption',
        'verification_source': 'host_verified_complete_terminal_namespace',
        'verification_method': 'registered_public_upstream_namespace',
        'proof_origin': 'client_complete_sse_terminal',
        'canonical_public_model': payload['model'], 'event_public_model': event['model'],
        'supplier_model_verbatim': candidate['model'], 'upstream_model': binding['supplier_model_verbatim'],
        'response_id': candidate['id'], 'raw_response_object_sha256': sha(candidate),
        'mapping_registry_id': registry['registry_id'],
        'mapping_registry_sha256': MODEL_NAMESPACE_REGISTRY_SHA256,
        'mapping_evidence': [{'kind': s['kind'], 'path': s['path'], 'sha256': s['sha256']} for s in binding['sources']],
        'no_model_field_rewrite': True}


def _transport_require_verified_identity(response, requested_public_model):
    identity = response.get('verified_model_identity') if isinstance(response, dict) else None
    if not isinstance(identity, dict):
        raise ModelNamespaceIdentityError('Client-generated or trusted-journal model identity required')
    if identity.get('verification_method') == 'exact_raw_response_model':
        expected = {'provider': 'adaption', 'verification_source': 'host_verified_complete_exact_response',
            'canonical_public_model': requested_public_model, 'supplier_model_verbatim': requested_public_model,
            'upstream_model': None, 'response_id': response.get('id'),
            'mapping_registry_id': None, 'mapping_registry_sha256': None, 'mapping_evidence': [],
            'no_model_field_rewrite': True}
        if any(identity.get(k) != v for k, v in expected.items()) or response.get('model') != requested_public_model or identity.get('proof_origin') not in ('client_complete_sse_terminal', 'client_complete_http_json'):
            raise ModelNamespaceIdentityError('Invalid exact-model identity')
        if identity.get('event_public_model') not in (None, requested_public_model):
            raise ModelNamespaceIdentityError('Conflicting exact-model public tag')
        return identity
    registry, binding = _transport_namespace_registration(requested_public_model)
    expected = {'provider': 'adaption', 'verification_source': 'host_verified_complete_terminal_namespace',
        'verification_method': 'registered_public_upstream_namespace', 'proof_origin': 'client_complete_sse_terminal',
        'canonical_public_model': requested_public_model, 'event_public_model': requested_public_model,
        'supplier_model_verbatim': binding['supplier_model_verbatim'], 'upstream_model': binding['supplier_model_verbatim'],
        'response_id': response.get('id'), 'mapping_registry_id': registry['registry_id'],
        'mapping_registry_sha256': MODEL_NAMESPACE_REGISTRY_SHA256, 'no_model_field_rewrite': True}
    if any(identity.get(k) != v for k, v in expected.items()) or response.get('model') != binding['supplier_model_verbatim']:
        raise ModelNamespaceIdentityError('A verified final public/upstream namespace identity is required')
    expected_sources = [{'kind': s['kind'], 'path': s['path'], 'sha256': s['sha256']} for s in binding['sources']]
    if identity.get('mapping_evidence') != expected_sources:
        raise ModelNamespaceIdentityError('Verified identity mapping evidence differs')
    return identity


def _transport_response_for_loop(raw_response, verified_identity):
    if not isinstance(raw_response, dict):
        return raw_response
    result = dict(raw_response)
    result.pop('verified_model_identity', None)  # Never trust a supplier field with this host-only name.
    if verified_identity is not None:
        result['verified_model_identity'] = verified_identity
    return result


def _transport_public_model_fields(value, key, organization_id, path='$'):
    found = []
    if isinstance(value, dict):
        for name, item in value.items():
            child = path + '.' + str(name)
            if name == 'model' and isinstance(item, str):
                found.append({'path': child, 'value': _transport_safe_text(item, key, organization_id)})
            if isinstance(item, (dict, list)):
                found.extend(_transport_public_model_fields(item, key, organization_id, child))
    elif isinstance(value, list):
        for index, item in enumerate(value):
            found.extend(_transport_public_model_fields(item, key, organization_id, path + '[' + str(index) + ']'))
    return found


def _transport_sse_remaining(client, deadline):
    remaining = deadline - time.monotonic()
    if remaining <= 0:
        raise OutcomeUnknown('SSE absolute request deadline exhausted; original request not repeated')
    sock = getattr(client.active, 'sock', None)
    if sock is not None:
        sock.settimeout(min(client.timeout, remaining))
    return remaining


def _transport_read_sse(client, response, row, identity, payload, deadline):
    raw = bytearray()
    buffer = bytearray()
    fields = {'data': [], 'event': '', 'id': None}
    counts = {}
    seen = {}
    events = duplicates = terminal_events = 0
    terminal = None
    terminal_kind = None
    observed_id = observed_model = None
    terminal_model_observed = None
    intermediate_model_observations = []
    terminal_response_observed = None
    eof = False
    last_checkpoint = 0
    last_error = None
    line_count = 0

    def checkpoint(phase):
        diagnostic = _transport_body_diagnostic(bytes(raw), client.key, client.org, eof)
        diagnostic['classification'] = 'event_stream'
        row.update(transport_phase=phase, requested_stream=True, actual_representation='sse',
                   http_body_diagnostic=diagnostic,
                   sse_diagnostic={'version': SSE_TRANSPORT_VERSION,
                     'bytes_observed': len(raw), 'raw_sha256': diagnostic['raw_sha256'],
                     'stream_eof_observed': eof, 'event_count': events, 'event_counts': dict(counts),
                     'duplicate_events': duplicates, 'terminal_events_observed': terminal_events,
                     'terminal_type': terminal_kind, 'response_id': observed_id, 'model': observed_model,
                     'terminal_model_observed': terminal_model_observed,
                     'intermediate_model_observations': list(intermediate_model_observations),
                     'terminal_response_observed': terminal_response_observed,
                     'model_confirmation_policy': 'Only exact registered final public/upstream namespaces with pinned evidence confirm deployment; raw supplier model remains unchanged.',
                     'verified_model_identity': row.get('verified_model_identity'),
                     'terminal_status': terminal.get('status') if terminal else None,
                     'usage': terminal.get('usage') if terminal else None,
                     'last_error_event_safe': last_error,
                     'stop_policy': 'Only full terminal response; check already-buffered frames, then close without waiting for stream EOF.'})
        client.journal.save(identity, row)

    def fail(category, message):
        row['error_category'] = category
        raise OutcomeUnknown(message + '; original request not repeated')

    def dispatch():
        nonlocal events, duplicates, terminal_events, terminal, terminal_kind
        nonlocal observed_id, observed_model, last_error, terminal_model_observed, terminal_response_observed
        if not fields['data']:
            fields.update(data=[], event='', id=None)
            return
        data = '\n'.join(fields['data'])
        event_name = fields['event']
        event_id = fields['id']
        fields.update(data=[], event='', id=None)
        events += 1
        if data.strip() == '[DONE]':
            counts['[DONE]'] = counts.get('[DONE]', 0) + 1
            return
        try:
            value = json.loads(data)
        except ValueError:
            fail('http_response_sse_event_json_decode', 'SSE data event is not JSON')
        if not isinstance(value, dict):
            fail('http_response_sse_event_shape', 'SSE data event must be an object')
        kind = value.get('type') or event_name
        if not isinstance(kind, str) or not kind:
            fail('http_response_sse_event_shape', 'SSE semantic event type missing')
        if event_name and event_name != 'message' and value.get('type') and event_name != value['type']:
            fail('http_response_sse_event_type_mismatch', 'SSE event and data types differ')
        counts[kind] = counts.get(kind, 0) + 1
        sequence = value.get('sequence_number')
        if sequence is not None and type(sequence) is not int:
            fail('http_response_sse_event_shape', 'SSE sequence number must be an integer')
        dedup_key = ('id', event_id) if event_id else ('sequence', sequence) if sequence is not None else None
        canonical = json.dumps(value, sort_keys=True, ensure_ascii=False)
        if dedup_key is not None and dedup_key in seen:
            if seen[dedup_key] != canonical:
                fail('http_response_sse_duplicate_conflict', 'Conflicting duplicate SSE event')
            duplicates += 1
            return
        if dedup_key is not None:
            seen[dedup_key] = canonical
        event_response_id = value.get('response_id')
        if event_response_id is not None:
            if not isinstance(event_response_id, str) or not event_response_id or (observed_id and observed_id != event_response_id):
                fail('http_response_sse_response_id_mismatch', 'SSE event response identity changed')
            observed_id = event_response_id
        candidate = value.get('response')
        if kind in ('response.completed', 'response.incomplete', 'response.failed'):
            terminal_response_observed = {
                'event_type': kind, 'event_object_sha256': sha(value),
                'response_object_sha256': sha(candidate),
                'id': _transport_safe_text(candidate.get('id'), client.key, client.org) if isinstance(candidate, dict) else None,
                'model': _transport_safe_text(candidate.get('model'), client.key, client.org) if isinstance(candidate, dict) else None,
                'status': candidate.get('status') if isinstance(candidate, dict) else None,
                'usage_safe_json': _transport_safe_text(json.dumps(candidate.get('usage'), sort_keys=True, ensure_ascii=False), client.key, client.org) if isinstance(candidate, dict) else None,
                'public_model_fields': _transport_public_model_fields(value, client.key, client.org),
                'safe_response_json': _transport_safe_text(json.dumps(candidate, sort_keys=True, ensure_ascii=False), client.key, client.org),
                'safe_response_truncated': False, 'accepted_as_terminal': False}
        if isinstance(candidate, dict):
            rid, model = candidate.get('id'), candidate.get('model')
            if rid is not None:
                if not isinstance(rid, str) or not rid or (observed_id and observed_id != rid):
                    fail('http_response_sse_response_id_mismatch', 'SSE response identity changed')
                observed_id = rid
            if model is not None:
                if kind in ('response.completed', 'response.incomplete', 'response.failed'):
                    terminal_model_observed = model
                else:
                    intermediate_model_observations.append({'event': kind,
                        'model': _transport_safe_text(model, client.key, client.org),
                        'outer_model': _transport_safe_text(value.get('model'), client.key, client.org) if value.get('model') is not None else None,
                        'public_model_fields': _transport_public_model_fields(value, client.key, client.org)})
        if kind in ('error', 'response.error'):
            last_error = _transport_safe_text(canonical, client.key, client.org, 2048)
            fail('http_response_sse_error_event', 'SSE error event without full terminal response')
        if kind in ('response.completed', 'response.incomplete', 'response.failed'):
            terminal_events += 1
            expected = kind.split('.', 1)[1]
            # The gateway can label a max-output terminal event completed while
            # the full response truthfully says incomplete. Keep both labels;
            # never execute tools or claim completion from such a response.
            capped_terminal = (kind == 'response.completed' and isinstance(candidate, dict)
                and candidate.get('status') == 'incomplete'
                and (candidate.get('incomplete_details') or {}).get('reason') == 'max_output_tokens'
                and candidate.get('max_output_tokens') == payload.get('max_output_tokens')
                and (candidate.get('usage') or {}).get('output_tokens') == payload.get('max_output_tokens'))
            if not isinstance(candidate, dict) or candidate.get('object') != 'response' or (candidate.get('status') != expected and not capped_terminal):
                fail('http_response_sse_terminal_shape', 'SSE terminal event lacks matching full response')
            if not isinstance(candidate.get('id'), str) or not candidate['id'] or not isinstance(candidate.get('model'), str) or not isinstance(candidate.get('output'), list):
                fail('http_response_sse_terminal_shape', 'SSE terminal response missing id/model/output')
            try:
                verified = _transport_verified_namespace_identity(payload, value, candidate, observed_id)
            except (ModelNamespaceIdentityError, ValueError, KeyError, OSError):
                fail('http_response_sse_model_identity', 'Final public/upstream namespace identity is not exactly registered with valid evidence')
            row['verified_model_identity'] = verified
            terminal_response_observed['verified_model_identity'] = verified
            observed_model = verified['canonical_public_model']
            if terminal is not None:
                if candidate != terminal:
                    fail('http_response_sse_duplicate_terminal_conflict', 'Conflicting SSE terminal responses')
                duplicates += 1
            else:
                terminal, terminal_kind = candidate, kind
                terminal_response_observed['accepted_as_terminal'] = True
        elif terminal is not None:
            fail('http_response_sse_event_after_terminal', 'Unexpected event after SSE terminal')

    def line_received(line):
        nonlocal line_count
        if line.endswith(b'\r'):
            line = line[:-1]
        try:
            text = line.decode('utf-8-sig' if line_count == 0 else 'utf-8')
            line_count += 1
        except UnicodeDecodeError:
            fail('http_response_sse_event_encoding', 'SSE event is not UTF8')
        if not text:
            dispatch()
            return
        if text.startswith(':'):
            return
        name, separator, value = text.partition(':')
        if separator and value.startswith(' '):
            value = value[1:]
        if name == 'data':
            fields['data'].append(value)
        elif name == 'event':
            fields['event'] = value
        elif name == 'id' and '\x00' not in value:
            fields['id'] = value
        # retry and unknown fields are ignored; no transport retry is enabled.

    try:
        checkpoint('http_response_sse_read')
        while terminal is None:
            chunk = None
            try:
                _transport_sse_remaining(client, deadline)
                chunk = response.readline(min(65536, MAX_RESPONSE - len(raw) + 1))
                _transport_sse_remaining(client, deadline)
            except BaseException as exc:
                if isinstance(chunk, bytes):
                    raw.extend(chunk)
                partial = getattr(exc, 'partial', None)
                if isinstance(partial, bytes):
                    raw.extend(partial)
                if 'error_category' not in row:
                    row['error_category'] = 'http_response_sse_deadline' if time.monotonic() >= deadline else 'http_response_sse_read'
                raise
            if not isinstance(chunk, bytes):
                fail('http_response_sse_read_shape', 'SSE reader returned non-bytes')
            if not chunk:
                eof = True
                break
            raw.extend(chunk)
            if len(raw) > MAX_RESPONSE:
                fail('http_response_sse_body_size', 'SSE response exceeds original response byte limit')
            buffer.extend(chunk)
            while b'\n' in buffer:
                line, rest = buffer.split(b'\n', 1)
                buffer[:] = rest
                line_received(bytes(line))
            if len(raw) - last_checkpoint >= 65536:
                checkpoint('http_response_sse_read')
                last_checkpoint = len(raw)
        if terminal is None:
            fail('http_response_sse_missing_terminal', 'SSE stream ended without a complete terminal frame')
        _transport_sse_remaining(client, deadline)
        checkpoint('http_response_sse_terminal_ready')
        return terminal
    except BaseException:
        checkpoint('http_response_sse_failed')
        raise



def _transport_rate_limit_headers(response, key, organization_id):
    names=('retry-after','date','ratelimit','ratelimit-policy','ratelimit-limit','ratelimit-remaining','ratelimit-reset',
           'x-ratelimit-limit','x-ratelimit-remaining','x-ratelimit-reset','x-ratelimit-limit-requests',
           'x-ratelimit-remaining-requests','x-ratelimit-reset-requests','x-ratelimit-limit-tokens',
           'x-ratelimit-remaining-tokens','x-ratelimit-reset-tokens')
    pairs=response.getheaders() if hasattr(response,'getheaders') else [(name,response.getheader(name))for name in names]
    result={}
    for name,value in pairs:
        name=name.lower()
        if value is not None and (name in names or name.startswith(('ratelimit-','x-ratelimit-','x-rate-limit-','x-concurrency-'))):
            result[name]=_transport_safe_text(str(value),key,organization_id,512)
    return result


class Client:
    def __init__(self, journal, key=None, organization_id=None, timeout=1800, base_url=API_BASE,
                 connection_factory=http.client.HTTPSConnection):
        self.journal = journal
        if key is None:
            key, org, self.credential_source = credentials()
            organization_id = organization_id or org
        else:
            self.credential_source = 'injected'
        self.key, self.org = key, organization_id
        self.timeout = timeout
        self.base = urllib.parse.urlsplit(base_url)
        if self.base.scheme != 'https' or self.base.hostname != 'api.prod.adaptionlabs.ai' or self.base.query:
            raise ValueError('Adaption requests require the official HTTPS API origin')
        self.connection_factory = connection_factory
        self.cancel_event = threading.Event()
        self.active = None
        self.latest_inference_request_identity = None
        self.remote = {'status': 'not_started', 'response_id': None}

    def cancel(self):
        self.cancel_event.set()
        if self.active is not None:
            self.active.close()
        # Closing a local socket does not cancel server-side inference.
        return {'local_cancel_requested': True, 'local_connection_closed': True,
                'remote_request_status': self.remote.get('status') if self.remote.get('status') in
                ('completed', 'failed', 'cancelled', 'incomplete') else 'unknown',
                'remote_cancel_attempted': False, 'response_id': self.remote.get('response_id'),
                'latest_inference_request_identity': self.latest_inference_request_identity}

    def request(self, method, suffix, identity, payload=None):
        fingerprint = sha({'method': method, 'suffix': suffix, 'payload': payload})
        with self.journal.lock(identity):
            prior = self.journal.load(identity)
            if prior:
                if prior['fingerprint'] != fingerprint:
                    raise ValueError('Request identity has different immutable input')
                if prior['state'] == 'received':
                    if not 200 <= prior['http_status'] < 300:
                        raise RuntimeError('Adaption HTTP ' + str(prior['http_status']) + '; original error receipt retained')
                    return _transport_response_for_loop(prior['response'], prior.get('verified_model_identity'))
                if prior['state']=='not_sent_budget_denied':raise RuntimeError('Original POST not sent: budget denied; immutable ID retained')
                if method == 'POST':
                    raise OutcomeUnknown('Original POST outcome unknown; query known response id, never resubmit')
            if self.cancel_event.is_set():
                raise LocalCancelled('Local cancellation already requested')
            row = {'version': VERSION, 'identity': identity, 'fingerprint': fingerprint,
                   'state': 'reserved', 'method': method, 'path': self.base.path + suffix,
                   'request': payload, 'created': time.time(), 'automatic_retries': 0,
                   'credential_source': self.credential_source, 'organization_id': self.org,
                   'remote_request_status': 'unknown' if method == 'POST' else 'not_inference'}
            reservation = None
            if method == 'POST':
                try:
                    import batch_budget_guard
                    grant_path = Path(os.environ['YOHAKU_BATCH_SPEND_GRANT'])
                    if sha(grant_path.read_bytes()) != os.environ['YOHAKU_BATCH_SPEND_GRANT_SHA256']:raise ValueError('Approved spending grant changed before POST')
                    reservation = batch_budget_guard.reserve_post(os.environ['YOHAKU_BATCH_ID'], identity, payload)
                except Exception as exc:
                    row.update(state='not_sent_budget_denied',network_dispatched=False,error_type=type(exc).__name__,remote_request_status='not_sent')
                    self.journal.save(identity,row)
                    raise
            self.journal.save(identity,row)
            if method == 'POST':
                # A prior round's completion is not the state of this new request.
                self.remote = {'status': 'unknown', 'response_id': None}
                self.latest_inference_request_identity = identity
            body = json.dumps(payload, ensure_ascii=False).encode() if payload is not None else None
            headers = {'Authorization': 'Bearer ' + self.key, 'Content-Type': 'application/json',
                       'User-Agent': VERSION, 'Connection': 'close'}
            if self.org:
                headers['X-Organization-Id'] = self.org
            request_deadline = time.monotonic() + self.timeout
            if payload is not None and payload.get('stream') is True:
                headers['Accept'] = 'text/event-stream'
            conn = self.connection_factory(self.base.hostname, timeout=self.timeout)
            self.active = conn
            try:
                conn.request(method, self.base.path + suffix, body=body, headers=headers)
                if payload is not None and payload.get('stream') is True:
                    _transport_sse_remaining(self, request_deadline)
                response = conn.getresponse()
                row.update(transport_observability_version=HTTP_DIAGNOSTICS_VERSION,
                           transport_phase='http_response_headers_received',
                           http_status=response.status,
                           rate_header_capture_version='retry-after-capture/1',
                           rate_limit_headers=_transport_rate_limit_headers(response,self.key,self.org),
                           http_content_type=_transport_safe_text(response.getheader('content-type') or '', self.key, self.org, 512),
                           server_request_id=_transport_safe_text(response.getheader('x-request-id') or '', self.key, self.org, 512))
                self.journal.save(identity, row)
                row.update(requested_stream=bool(payload is not None and payload.get('stream') is True))
                content_type = (response.getheader('content-type') or '').split(';', 1)[0].strip().lower()
                if row['requested_stream'] and 200 <= response.status < 300 and content_type == 'text/event-stream':
                    value = _transport_read_sse(self, response, row, identity, payload, request_deadline)
                else:
                    row['actual_representation'] = 'non_sse'
                    try:
                        if row['requested_stream']:
                            _transport_sse_remaining(self, request_deadline)
                        raw = response.read(MAX_RESPONSE + 1)
                        row['http_body_diagnostic'] = _transport_body_diagnostic(raw, self.key, self.org, len(raw) <= MAX_RESPONSE)
                        if row['requested_stream']:
                            _transport_sse_remaining(self, request_deadline)
                    except BaseException as read_exc:
                        partial = getattr(read_exc, 'partial', None)
                        row.update(transport_phase='http_response_body_read',
                                   error_category='http_response_sse_deadline' if row['requested_stream'] and time.monotonic() >= request_deadline else 'http_response_body_read',
                                   body_read_exception_type=type(read_exc).__name__)
                        if isinstance(partial, bytes):
                            row['http_body_diagnostic'] = _transport_body_diagnostic(partial, self.key, self.org, False)
                        self.journal.save(identity, row)
                        raise
                    row.update(transport_phase='http_response_json_decode',
                               http_body_diagnostic=_transport_body_diagnostic(raw, self.key, self.org, len(raw) <= MAX_RESPONSE))
                    self.journal.save(identity, row)
                    if len(raw) > MAX_RESPONSE:
                        row.update(error_category='http_response_body_size')
                        self.journal.save(identity, row)
                        raise OutcomeUnknown('Response too large; original request not repeated')
                    try:
                        value = json.loads(raw)
                    except ValueError as decode_exc:
                        row.update(error_category='http_response_json_decode',
                                   http_json_decode_exception_type=type(decode_exc).__name__)
                        self.journal.save(identity, row)
                        raise OutcomeUnknown('Non-JSON HTTP response; original request not repeated') from None
                    row.update(transport_phase='http_response_json_decoded', http_json_envelope_type=type(value).__name__)
                    row['actual_representation'] = 'json'
                    if row['requested_stream'] and 200 <= response.status < 300 and not (content_type == 'application/json' or content_type.endswith('+json')):
                        row.update(error_category='http_response_sse_content_type')
                        self.journal.save(identity, row)
                        raise OutcomeUnknown('Unexpected non-SSE response content type; original request not repeated')
                # Host recovery: non-2xx is an HTTP error, not a model response.
                # The existing worst-case POST reservation remains untouched.
                if not 200 <= response.status < 300:
                    row.update(state='received', error_category='http_response_status',
                               remote_request_status='unknown', finished=time.time())
                    self.journal.save(identity, row)
                    raise RuntimeError('Adaption HTTP ' + str(response.status) + '; see redacted receipt ' + identity)
                if method == 'POST' and row.get('verified_model_identity') is None:
                    try:
                        row['verified_model_identity'] = _transport_verified_exact_identity(payload, value, 'client_complete_http_json')
                    except ModelNamespaceIdentityError:
                        row.update(response_observed_unverified=value, error_category='http_response_model_identity')
                        self.journal.save(identity, row)
                        raise OutcomeUnknown('Complete response lacks an exact or verified namespaced model identity; original request not repeated')
                if reservation is not None and isinstance(value,dict):
                    accounting_response = _transport_response_for_loop(value, row.get('verified_model_identity'))
                    batch_budget_guard.reconcile_post(reservation, accounting_response)
                row.update(state='received', http_status=response.status, response=value,
                           server_request_id=_transport_safe_text(response.getheader('x-request-id') or '', self.key, self.org, 512), finished=time.time())
                if isinstance(value, dict) and value.get('object') == 'response':
                    self.remote = {'status': value.get('status', 'unknown'), 'response_id': value.get('id')}
                    row['remote_request_status'] = self.remote['status']
                self.journal.save(identity, row)
                if not 200 <= response.status < 300:
                    raise RuntimeError('Adaption HTTP ' + str(response.status) + '; see redacted receipt ' + identity)
                return _transport_response_for_loop(value, row.get('verified_model_identity'))
            except BaseException as exc:
                if row['state'] != 'received':
                    row.update(state='local_cancelled' if self.cancel_event.is_set() else 'outcome_unknown',
                               local_exception_type=type(exc).__name__, finished=time.time(),
                               remote_request_status='unknown', remote_cancel_attempted=False)
                    self.journal.save(identity, row)
                raise
            finally:
                conn.close(); self.active = None

    def models(self, identity='models-hosted'):
        return self.request('GET', '/models?type=hosted', identity)

    def create(self, payload, identity):
        return self.request('POST', '/responses', identity, payload)

    def retrieve(self, response_id, identity):
        if (not isinstance(response_id, str) or not response_id.startswith('resp_') or
                not 5 < len(response_id) <= 8192 or
                any(unicodedata.category(c) == 'Cc' for c in response_id)):
            raise ValueError('Invalid response id')
        # Supplier IDs are opaque and may contain padding or path delimiters.
        # The resp_ prefix excludes standalone dot segments, and encoding the
        # entire component keeps '/', '?', '#', '%' and '=' out of URL structure.
        component = urllib.parse.quote(response_id, safe='')
        return self.request('GET', '/responses/' + component, identity)


def function_definitions(definitions):
    return [{'type': 'function', 'name': 'yohaku_' + d['name'],
             'description': d['description'], 'parameters': d['parameters']} for d in definitions]


def output_text(response):
    return ''.join(c.get('text', '') for item in response.get('output', []) if item.get('type') == 'message'
                   for c in item.get('content', []) if c.get('type') == 'output_text')


def inherit_profile(parent, requested=None):
    requested = requested or {}
    chosen = json.loads(json.dumps(parent)); chosen.update(requested)
    if chosen.get('provider') != 'adaption' or chosen.get('model') != parent['model']:
        raise ValueError('Adaption descendants must inherit provider/model; no silent fallback')
    if chosen.get('reasoning_effort') != 'none':
        raise ValueError('Adaption reasoning effort is not exposed by this API contract; use none')
    return chosen


# four-lines-completed-turn-checkpoint/1
def _completed_turn_checkpoint(loop, history, response, round_no):
    node = getattr(loop.tools, 'node_context', {})
    value = {'format': 'adaption-completed-turn-checkpoint/1',
        'task_id': loop.identity, 'node_id': node.get('node_id'),
        'generation': node.get('generation'), 'profile': node.get('profile'),
        'session_id': loop.session_id, 'native_tools_sha256': sha(function_definitions(loop.tools.definitions)),
        'completed_round': round_no, 'response_id': response.get('id'),
        'history_sha256': sha(history),
        'continuation': {'history': history, 'last_response_id': response.get('id'), 'store': loop.store, 'model': loop.model},
        'boundary': 'Complete supplier terminal and all returned local tool outputs; never includes an unknown later POST.'}
    loop.client.journal.save(loop.identity + ':checkpoint', value)

def _completed_checkpoint_ancestor(previous_task, node, session_id=None):
    expected_session = session_id or node.get('session_id')
    expected_profile = node.get('profile')
    if not expected_session:
        raise OutcomeUnknown('No confirmed original session for checkpoint continuity')
    base = Path.home() / '.local/state/yohaku'
    current = previous_task
    seen = set()
    skipped = []
    while current and current not in seen:
        seen.add(current)
        if not re.fullmatch(r'model-task-[a-f0-9]+', current):
            raise ValueError('Invalid local checkpoint task identity')
        path = base / 'adaption-tasks' / current / (current + ':checkpoint.json')
        if path.is_file() and not path.is_symlink():
            value = json.loads(path.read_bytes())
            if (value.get('task_id') != current or value.get('node_id') != node.get('node_id')
                    or value.get('generation') != node.get('generation') or value.get('profile') != expected_profile
                    or value.get('session_id') != expected_session
                    or value.get('history_sha256') != sha(value.get('continuation', {}).get('history'))):
                raise ValueError('Completed checkpoint ancestor identity/profile/history mismatch')
            if any(x['restored_history_sha256'] != value['history_sha256'] for x in skipped):
                raise ValueError('No-tool carry-forward does not match the actual restored history')
            return value, skipped
        record_path = base / 'model-tasks' / (current + '.json')
        if not record_path.is_file() or record_path.is_symlink():
            raise OutcomeUnknown('No durable previous task record for checkpoint continuity')
        record = json.loads(record_path.read_bytes())
        context = record.get('node_context') or {}
        if (record.get('task_id') != current or record.get('model_phase') != 'finished'
                or record.get('commit_confirmed') is not True
                or (record.get('receipt') or {}).get('native_cleanup_complete') is not True
                or context.get('node_id') != node.get('node_id') or context.get('generation') != node.get('generation')
                or context.get('profile') != expected_profile or context.get('session_id') != expected_session):
            raise OutcomeUnknown('Checkpoint carry-forward requires exact confirmed terminal task/context/cleanup')
        trace = base / 'model-tasks' / (current + '.node-events.jsonl')
        events = [json.loads(line) for line in trace.read_text().splitlines()]
        if any(e.get('event') in ('tool_call', 'tool_result') for e in events):
            raise OutcomeUnknown('Task has tool events without a completed checkpoint; no replay or skip')
        restores = [e for e in events if e.get('event') == 'completed_turn_continuation_restored']
        previous = context.get('adaption_previous_task_id')
        if len(restores) != 1 or restores[0].get('previous_task_id') != previous or restores[0].get('session_id') != expected_session:
            raise OutcomeUnknown('No actual restore receipt before the no-tool HTTP failure')
        rounds = list((base / 'adaption-tasks' / current).glob(current + ':round:*.json'))
        if len(rounds) != 1:
            raise OutcomeUnknown('Only a first-POST HTTP429 with no tool events can carry its inherited checkpoint')
        row = json.loads(rounds[0].read_bytes())
        if (row.get('identity') != current + ':round:0' or row.get('finished') is None
                or (row.get('request') or {}).get('model') != expected_profile['model']):
            raise OutcomeUnknown('Missing exact first-POST terminal receipt; no checkpoint skip')
        response = row.get('response') or {}
        verified = row.get('verified_model_identity') or {}
        known_throttle = row.get('http_status') == 429
        known_no_action_cap = (row.get('http_status') == 200 and row.get('state') == 'received'
            and response.get('object') == 'response' and response.get('status') == 'incomplete'
            and (response.get('incomplete_details') or {}).get('reason') == 'max_output_tokens'
            and (response.get('usage') or {}).get('output_tokens') == (row.get('request') or {}).get('max_output_tokens')
            and response.get('model') == verified.get('supplier_model_verbatim')
            and verified.get('canonical_public_model') == expected_profile['model']
            and verified.get('raw_response_object_sha256') == sha(response)
            and not any(x.get('type') == 'function_call' for x in response.get('output', []) if isinstance(x, dict)))
        if not (known_throttle or known_no_action_cap):
            raise OutcomeUnknown('No confirmed no-action HTTP429/output-limit terminal; no checkpoint skip')
        skipped.append({'task_id':current,'journal_path':str(rounds[0]),'journal_sha256':sha(rounds[0].read_bytes()),
                        'restored_history_sha256':restores[0].get('history_sha256'),'http_status':row.get('http_status'),
                        'terminal_kind':'confirmed_output_limit' if known_no_action_cap else 'http429','tool_events':0})
        current = previous
    raise OutcomeUnknown('Checkpoint ancestry is absent or cyclic; no inferred continuation')

def _completed_prior_result(prior, previous_task, node, tools):
    result = prior.get('final_result') or {}
    if isinstance(result.get('continuation', {}).get('history'), list):
        return result
    value, skipped = _completed_checkpoint_ancestor(previous_task, node)
    if value.get('native_tools_sha256') != sha(function_definitions(tools.definitions)):
        raise ValueError('Completed checkpoint tool catalog mismatch')
    if getattr(tools, 'node_event', None):
        tools.node_event({'event':'completed_turn_continuation_restored','previous_task_id':previous_task,
            'checkpoint_task_id':value['task_id'],'skipped_no_tool_terminal_tasks':skipped,
            'session_id':value['session_id'],'history_sha256':value['history_sha256'],
            'completed_round':value['completed_round'],'unknown_later_POSTs_replayed':False})
    return value

def _checkpoint_native_signature(task_id, node, session_id):
    direct = Path.home() / '.local/state/yohaku/adaption-tasks' / task_id / (task_id + ':checkpoint.json')
    if not direct.is_file():
        record = Path.home() / '.local/state/yohaku/model-tasks' / (task_id + '.json')
        if not record.is_file() or not (json.loads(record.read_bytes()).get('node_context') or {}).get('adaption_previous_task_id'):
            return {}
    value, skipped = _completed_checkpoint_ancestor(task_id, node, session_id or node.get('session_id'))
    if skipped:
        return dict(value, checkpoint_ancestor_carry_forward=True, skipped_no_tool_http429_tasks=skipped)
    return value


class Loop:
    def __init__(self, client, tools, model, request_id, *, seconds=None, max_tools=None,
                 max_output_tokens=32768, emit=lambda event: None, store=False):
        self.client, self.tools, self.model, self.identity = client, tools, model, request_id
        self.seconds, self.max_tools = seconds, max_tools
        self.max_output_tokens, self.emit, self.store = max_output_tokens, emit, store
        self.session_id = 'adaption-context-' + sha(self.identity)[:32]

    def run(self, prompt, context=None, history=None, instructions='Use only the provided Yohaku world tools.'):
        begin = time.monotonic()
        context = context or {}
        definitions = function_definitions(self.tools.definitions)
        allowed = {d['name'] for d in definitions}
        input_items = list(history or []) + [{'role': 'user', 'content': prompt}]
        full_history = list(input_items)
        usage, receipts, local_receipts, calls_total = [], [], [], 0
        previous = None
        self.emit({'event': 'requested_model_settings', 'provider': 'adaption', 'model': self.model,
                   'reasoning_effort': 'none', 'selection_source': 'explicit immutable node profile',
                   'scope': 'isolated Adaption candidate only'})
        self.emit({'event': 'turn_settings', 'version': VERSION, 'seconds': self.seconds,
                   'max_tool_calls': self.max_tools, 'automatic_retries': 0,
                   'context_mode': 'stored_response_with_local_archive' if self.store else 'full_local_history',
                   'remote_cancellation': 'unknown until supplier receipt; no cancel endpoint used'})
        round_no=-1
        while True:
            round_no+=1
            remaining = None if self.seconds is None else self.seconds - (time.monotonic() - begin)
            if remaining is not None and remaining <= 0:
                raise BudgetExceeded('Adaption local wall-clock budget exhausted')
            self.client.timeout = 1800 if remaining is None else min(1800, remaining)
            payload = {'model': self.model, 'input': input_items, 'tools': definitions,
                       'instructions': instructions, 'stream': True, 'store': self.store,
                       'background': False, 'max_output_tokens': self.max_output_tokens}
            if previous:
                payload['previous_response_id'] = previous
            identity = self.identity + ':round:' + str(round_no)
            response = self.client.create(payload, identity)
            verified_identity = _transport_require_verified_identity(response, self.model)
            self.emit({'event': 'model_settings', 'provider': 'adaption',
                       'model': verified_identity['canonical_public_model'],
                       'canonical_public_model': verified_identity['canonical_public_model'],
                       'supplier_model_verbatim': response.get('model'),
                       'upstream_model': verified_identity['upstream_model'],
                       'verified_model_identity': verified_identity,
                       'mapping_evidence': verified_identity['mapping_evidence'],
                       'reasoning_effort': 'none',
                       'selection_source': 'complete supplier terminal plus host-pinned namespace evidence',
                       'response_id': response.get('id'), 'confirmed': True,
                       'scope': 'isolated Adaption candidate only'})
            if response.get('status') not in ('completed', 'incomplete'):
                raise RuntimeError('Adaption response not completed: ' + str(response.get('status')))
            output = response.get('output', [])
            if not isinstance(output, list):
                raise ValueError('Invalid response output')
            function_calls = [item for item in output if item.get('type') == 'function_call']
            ids = [c.get('call_id') for c in function_calls]
            if any(not isinstance(i, str) or not i for i in ids) or len(set(ids)) != len(ids):
                raise ValueError('Invalid/duplicate tool call IDs; no tools executed')
            if any(c.get('name') not in allowed for c in function_calls):
                raise ValueError('Unadvertised tool requested; no tools executed')
            if payload.get('stream') and function_calls and (response.get('status') != 'completed' or any(c.get('status') not in (None, 'completed') for c in function_calls)):
                raise RuntimeError('Streaming response has partial/incomplete tool calls; no tools executed')
            parsed_arguments = [json.loads(c.get('arguments') or '{}') for c in function_calls]
            if any(not isinstance(a, dict) for a in parsed_arguments):
                raise ValueError('Tool arguments must be JSON objects; no tools executed')
            if self.max_tools is not None and calls_total + len(function_calls) > self.max_tools:
                raise BudgetExceeded('Adaption tool-call budget exhausted before execution')
            full_history.extend(output)
            usage.append(response.get('usage'))
            receipts.append({'identity': identity, 'id': response.get('id'), 'status': response.get('status'),
                             'model': response.get('model'), 'usage': response.get('usage'),
                             'previous_response_id': previous, 'input_sha256': sha(input_items)})
            self.emit({'event': 'node_session', 'thread_id': self.session_id})
            self.emit({'event': 'node_usage', 'usage': {'source': 'Adaption Responses receipts', 'rounds': usage},
                       'elapsed_seconds': time.monotonic() - begin})
            if not function_calls:
                if response['status'] != 'completed':
                    raise RuntimeError('Incomplete response is not successful completion')
                result = {'text': output_text(response), 'session_id': self.session_id,
                          'native_tools_sha256': sha(definitions),
                          'local_tool_receipts': local_receipts,
                          'usage': {'source': 'Adaption Responses receipts', 'rounds': usage},
                          'elapsed_seconds': time.monotonic() - begin,
                          'provider_evidence': {'provider': 'adaption', 'requests': receipts,
                                                'local_cleanup': 'HTTP connections closed',
                                                'remote_request_status': response['status']},
                          'submitted_context': {'policy': 'scoped_yohaku_tools', 'context': context,
                                                'history_injected': bool(history)},
                          'continuation': {'history': full_history, 'last_response_id': response.get('id'),
                                           'store': self.store, 'model': self.model}}
                self.client.journal.save(self.identity + ':result', result)
                return result
            results = []
            for c, args in zip(function_calls, parsed_arguments):
                if self.client.cancel_event.is_set():
                    raise LocalCancelled('Local cancellation before tool execution')
                if self.seconds is not None and time.monotonic() - begin >= self.seconds:
                    raise BudgetExceeded('Adaption wall-clock budget exhausted before tool execution')
                tool_identity = self.identity + ':tool:' + sha([response.get('id'), c['call_id']])[:32]
                fingerprint = sha([c['name'], args])
                with self.client.journal.lock(tool_identity):
                    saved = self.client.journal.load(tool_identity)
                    executed_now = saved is None
                    if saved:
                        if saved['fingerprint'] != fingerprint:
                            raise ValueError('Tool call identity input conflict')
                        if saved['state'] != 'returned':
                            raise OutcomeUnknown('Local tool outcome unknown; never replay the operation')
                        tool_result = saved['result']
                    else:
                        self.client.journal.save(tool_identity, {'state': 'reserved', 'fingerprint': fingerprint,
                                                                 'name': c['name'], 'arguments': args,
                                                                 'call_id': c['call_id'], 'response_id': response.get('id')})
                        self.emit({'event': 'tool_call', 'tool': c['name'], 'arguments': args})
                        tool_result = self.tools.call(c['name'].removeprefix('yohaku_'), args)
                        self.client.journal.save(tool_identity, {'state': 'returned', 'fingerprint': fingerprint,
                                                                 'name': c['name'], 'arguments': args,
                                                                 'result': tool_result, 'call_id': c['call_id'],
                                                                 'response_id': response.get('id')})
                        self.emit({'event': 'tool_result', 'tool': c['name'], 'arguments': args, 'result': tool_result})
                local_receipts.append({'identity': tool_identity, 'name': c['name'], 'arguments': args,
                                       'result': tool_result, 'executed_now': executed_now,
                                       'call_id': c['call_id'], 'response_id': response.get('id')})
                calls_total += 1
                results.append({'type': 'function_call_output', 'call_id': c['call_id'],
                                'output': json.dumps(tool_result, ensure_ascii=False)})
            full_history.extend(results)
            _completed_turn_checkpoint(self, full_history, response, round_no)
            if self.store:
                previous = response['id']; input_items = results
            else:
                input_items = list(full_history)
        raise BudgetExceeded('Adaption round guard exhausted')


def catalog():
    try:
        key, org, source = credentials()
    except ValueError:
        return {'provider': 'adaption', 'configured': False, 'models': [],
                'login': {'state': 'not_signed_in'}, 'reason': 'Missing Adaption API credential'}
    journal = Journal(Path.home() / '.local/state/yohaku/adaption-catalog')
    client = Client(journal, key=key, organization_id=org)
    raw = client.models('hosted-' + str(time.time_ns()))
    return {'provider': 'adaption', 'configured': True, 'login': {'state': 'api_key_present'},
            'models': [{**m, 'executable_reasoning_efforts': ['none'],
                        'inference_tested': False} for m in raw['data']],
            'catalog_is_inference_test': False, 'credential_source': source,
            'compatibility': {'backend': 'Adaption Responses direct function tools',
                              'context': 'local full history; optional stored previous_response_id',
                              'remote_cancel': 'unverified/no endpoint invoked'}}


def run(initial, tools, emit):
    """Candidate task-runner entry, preserving native scoped/group-wrapped tools."""
    config = initial['config']
    node = getattr(tools, 'node_context', {})
    chosen = node.get('profile') or config.get('adaption_profile')
    if not chosen or chosen.get('provider') != 'adaption':
        raise ValueError('An explicit Adaption node profile is required')
    inherit_profile(chosen)
    task_id = initial.get('context', {}).get('task_id')
    if not task_id:
        raise ValueError('A durable task identity is required')
    journal = Journal(Path.home() / '.local/state/yohaku/adaption-tasks' / task_id)
    client = Client(journal)
    if getattr(tools, 'node_event', None):
        original_emit = emit
        def emit(event):
            tools.node_event(event); original_emit(event)
    history = None
    continuation = node.get('continuation')
    if continuation and continuation.get('model') == chosen['model']:
        history = continuation.get('history')
    previous_task = node.get('adaption_previous_task_id')
    if previous_task:
        from yohaku_codex_tasks import read_record
        prior = read_record(previous_task)
        prior_node = prior.get('node_context', {})
        previous_result = _completed_prior_result(prior, previous_task, node, tools)
        if (prior_node.get('node_id') != node.get('node_id') or
                prior_node.get('generation') != node.get('generation') or
                any(prior_node.get('profile', {}).get(k) != node.get('profile', {}).get(k)
                    for k in ('provider', 'model', 'reasoning_effort')) or
                previous_result.get('native_tools_sha256') != sha(function_definitions(tools.definitions))):
            raise ValueError('Previous Adaption context identity/profile/catalog mismatch')
        history = previous_result.get('continuation', {}).get('history')
        if not isinstance(history, list):
            raise OutcomeUnknown('Previous Adaption continuation unavailable; no silent history loss')
    # AI-node context includes only its own compact/current task and scoped assets.
    prompt = json.dumps({'task': initial['prompt'], 'node': node,
                         'bootstrap': tools.bootstrap(initial.get('context', {}))}, ensure_ascii=False)
    loop = Loop(client, tools, chosen['model'], task_id, seconds=chosen['budget']['timeout_seconds'],
                max_tools=chosen['budget']['max_tool_calls'], emit=emit)
    if node.get('session_id'):
        loop.session_id = node['session_id']
    try:
        return loop.run(prompt, initial.get('context'), history,
                        instructions=getattr(tools, 'node_system', 'Use only provided Yohaku tools.'))
    finally:
        journal.save(task_id + ':local-cleanup', client.cancel())


def main():
    p = argparse.ArgumentParser(description=__doc__)
    p.add_argument('mode', choices=['models', 'text', 'status', 'tool-loop'])
    p.add_argument('--journal', required=True)
    p.add_argument('--request-id', required=True)
    p.add_argument('--model')
    p.add_argument('--prompt')
    p.add_argument('--response-id')
    p.add_argument('--candidate')
    p.add_argument('--world')
    p.add_argument('--read-path')
    a = p.parse_args()
    c = Client(Journal(a.journal))
    def interrupted(number, _frame):
        receipt = c.cancel()
        receipt['signal'] = number
        c.journal.save(a.request_id + ':cancel', receipt)
        raise LocalCancelled('Local CLI cancelled; supplier cancellation is not confirmed')
    signal.signal(signal.SIGINT, interrupted)
    signal.signal(signal.SIGTERM, interrupted)
    if a.mode == 'models':
        result = c.models(a.request_id)
    elif a.mode == 'status':
        result = c.retrieve(a.response_id, a.request_id)
    elif a.mode == 'text':
        if not a.model or not a.prompt:
            p.error('text requires --model and --prompt')
        result = c.create({'model': a.model, 'input': a.prompt, 'tools': [],
                           'stream': False, 'store': False, 'background': False,
                           'max_output_tokens': 256}, a.request_id)
    else:
        if not all((a.model, a.candidate, a.world, a.read_path)):
            p.error('tool-loop requires --model --candidate --world --read-path')
        candidate = Path(a.candidate).resolve()
        # Installed copy's root is candidate/agent; source copy's is workflow/adaption.
        if (candidate / 'agent/ADAPTION-INSTALL.json').is_file() is False:
            raise ValueError('Read-only probe requires an installed isolated candidate')
        sys.path.insert(0, str(candidate / 'agent'))
        import yohaku_agent as agent
        world = Path(a.world).resolve()
        path = Path(a.read_path)
        if path.is_absolute() or '..' in path.parts or not str(path):
            raise ValueError('Probe needs an existing world-relative path')
        workspace = agent.Workspace(world, candidate / 'bin/loom-linux-x86_64', candidate / 'bin/yohaku_workspace')
        native_tools = agent.Tools(workspace.execute, lambda: {'root': str(world)},
                                   {'read': True, 'write': False, 'run': False, 'delete': False, 'import': False}, native=True)
        class ReadOnlyProbe:
            definitions = [d for d in native_tools.definitions if d['name'] == 'read_text']
            calls = []
            def call(self, name, args):
                if name != 'read_text' or args != {'path': a.read_path}:
                    raise ValueError('Probe grants exactly one read-only path')
                value = native_tools.call(name, args)
                self.calls.append({'name': name, 'arguments': args, 'result': value})
                return value
        tools = ReadOnlyProbe()
        prompt = a.prompt or ('Call yohaku_read_text exactly once with path ' + a.read_path +
                               '. Then summarize the returned text in one sentence and cite its revision. Do not answer before reading.')
        events = []
        result = Loop(c, tools, a.model, a.request_id, seconds=120, max_tools=1,
                      max_output_tokens=512, emit=events.append).run(prompt)
        actual_tools = result['local_tool_receipts']
        result['readonly_tool_loop_verified'] = (len(actual_tools) == 1 and
            actual_tools[0]['result'].get('ok') is True and len(result['provider_evidence']['requests']) >= 2 and bool(result['text']))
        result['actual_tool_receipts'] = actual_tools
        c.journal.save(a.request_id + ':probe', {'result': result, 'events': events})
    print(json.dumps(redact(result), ensure_ascii=False))


if __name__ == '__main__':
    main()
````

## 输出与关联
```yaml
output: "research/workflow-runtime/agent/yohaku_adaption.py"
```
